% Computing Threefold Flips in Macaulay2
% Companion note to the FlipComputation package.
% Build:  pdflatex IMPLEMENTATION.tex   (twice, for the references)
\documentclass[11pt,reqno]{amsart}

\usepackage[T1]{fontenc}
\usepackage[utf8]{inputenc}
\usepackage{amsmath,amssymb,amsthm}
\usepackage{booktabs,longtable,array}
\usepackage{microtype}
\usepackage[margin=1in]{geometry}
\usepackage{xcolor}
\usepackage[colorlinks=true,linkcolor=blue!50!black,urlcolor=blue!50!black,
            citecolor=blue!50!black]{hyperref}

% Verbatim blocks hold Macaulay2 input and output.
\usepackage{fancyvrb}
\fvset{fontsize=\small,xleftmargin=1em}

% pandoc emits \tightlist for lists without blank lines between items.
\providecommand{\tightlist}{%
  \setlength{\itemsep}{0pt}\setlength{\parskip}{0pt}}

% Long \texttt identifiers cannot hyphenate; give the paragraph builder room.
\setlength{\emergencystretch}{3em}

% Notation.  These are the commands GitHub's markdown renderer refuses, which is
% why the markdown version of this note had to spell them with Unicode letters.
\newcommand{\PP}{\ensuremath{\mathbb{P}}}
\newcommand{\QQ}{\ensuremath{\mathbb{Q}}}
\newcommand{\RR}{\ensuremath{\mathbb{R}}}
\newcommand{\bbA}{\ensuremath{\mathbb{A}}}
\newcommand{\sO}{\ensuremath{\mathcal{O}}}

\DeclareMathOperator{\Proj}{Proj}
\DeclareMathOperator{\Spec}{Spec}
\DeclareMathOperator{\divisor}{div}
\DeclareMathOperator{\Bl}{Bl}
\DeclareMathOperator{\codim}{codim}
\DeclareMathOperator{\Exc}{Exc}

\theoremstyle{plain}
\newtheorem{theorem}{Theorem}[section]
\newtheorem{lemma}[theorem]{Lemma}
\newtheorem{proposition}[theorem]{Proposition}
\theoremstyle{definition}
\newtheorem{definition}[theorem]{Definition}
\newtheorem{remark}[theorem]{Remark}

\title{Computing Threefold Flips in Macaulay2}
\author{Takehiko Yasuda}
\date{August 2026}

\begin{document}

\begin{abstract}
We implement Algorithm 3 (Computing a flip) of the
author\textquotesingle s paper on the minimal model program in dimension
three, and run it on explicit threefolds. Given the coordinate ring of
the target \(X\) of a flipping contraction, the implementation produces
the flip as

\[Z = \Proj\, R[I^{(m)}t] \longrightarrow X,\]

together with the certificates of Lemma 7.2 --- that \(Z\) is normal and
that the exceptional locus has codimension at least two. The centrepiece
is a genuine three-dimensional flip of a non-\QQ{}-Gorenstein toric
singularity, computed from end to end and checked against toric geometry
derived independently of the algorithm. This note is not a formal
verification; it records a reproducible experiment and sets out
precisely which claims the computations do and do not support. The code
and this note were written essentially by an AI system, with edits by
the author; see Section~\ref{use-of-ai-in-the-implementation}.
\end{abstract}

\maketitle

\section{Purpose and background}\label{purpose-and-background}

Let \(f : Y \to X\) be a flipping contraction: a small projective
birational morphism of normal threefolds with \(-K_Y\) \(f\)-ample and
relative Picard rank one. Its \textbf{flip} is the small projective
birational \(g : Z \to X\) with \(K_Z\) \(g\)-ample. Existence is a
theorem; the question here is whether the flip can be \emph{computed}
from equations.

Algorithm 3 of the paper answers this by the classical construction

\[Z = \Proj_X \; \bigoplus_{k \geq 0} \; \sO_X(-kmE), \qquad E \sim -K_X,\]

which is a flip once \(m\) is chosen so that the algebra is finitely
generated and the result is normal and small. The algorithm searches for
such an \(m\) and verifies the two conditions rather than assuming them.

The input is the homogeneous coordinate ring \(R\) of \(X\), so that
\(X = \Proj\, R\) sits in a weighted projective space; the
output is a \textbf{B2M projection}, a bigraded variety
\(Z \subset \PP^{r-1} \times X\) together with its projection to \(X\)
(Definition 2.4 of the paper). Beyond the scope of the paper, the
implementation also accepts an affine base \(X = \Spec\, R\),
which is what makes computation at a single singularity cheap.

\section{What the algorithm
computes}\label{what-the-algorithm-computes}

Writing \(K_X = \sum_i n_i D_i\) with \(D_i = V(p_i)\) of height one,
the five steps are:

\begin{enumerate}
\def\labelenumi{\arabic{enumi}.}
\tightlist
\item
  compute \(K_X\);
\item
  choose a homogeneous \(s\) with \(-K_X + \divisor(s)\) effective;
\item
  form \(I = \sO_X(K_X - \divisor(s))\) and \(E\) with
  \(I = \sO_X(-E)\);
\item
  for \(m = 1!, 2!, 3!, \dots\);
\item
  build \(Z^m = \Proj\, R[I^{(m)}t]\) and test Lemma 7.2 ---
  return \(Z^m\) at the first \(m\) that passes.
\end{enumerate}

Steps 1--3 are \texttt{flipDivisorData}, step 5 is
\texttt{bigradedReesProjection} together with \texttt{isSmallProjection}
and \texttt{isS2Source}, and the loop is \texttt{computeFlip}.

\subsection{Where the implementation departs from the
paper}\label{where-the-implementation-departs-from-the-paper}

Four departures matter mathematically. The first is a correction; the
others are reformulations that leave the answer alone and change the
cost by orders of magnitude.
\href{https://github.com/takehikoyasuda/flip-computation/blob/main/docs/implementation-notes.md}{\texttt{docs/implementation-notes.md}}
gives the full argument and the measurements for each.

\textbf{Weights on the fibre, not levelled generators.} Presenting
\(Z^m\) inside \(X \times \PP^{r-1}\), with one fibre coordinate \(u_i\)
per generator \(f_i\) of \(I^{(m)}\), makes the Rees relations
bihomogeneous only if the \(f_i\) share a degree ---
\(u_i \mapsto f_i t\) forces \(\deg u_i = \deg f_i + \deg t\) to be
independent of \(i\). Replacing \(I^{(m)}\) by the ideal generated by
its degree-\(D\) part is the obvious repair and is what this package
originally did. \textbf{It is unsound over a weighted grading:} that
ideal is \(\sum_i R_{D-d_i} f_i\), and \(R_{D-d_i}\) can vanish,
dropping \(f_i\) outright. Grading the threefold below by
\(\ell = (3,2,1)\) gives ring degrees \(2,5,3,6,7\) and generators in
degrees 2 and 3, so \(D = 3\) needs \(R_1 = 0\) --- and the projection
is then reported as \emph{not small when it is small}. Enlarging \(D\)
does not help; the \(D\) that preserve the sheaf are sparse,
non-monotone, and sometimes absent. The implementation instead weights
the fibre direction, with \(d_0 = \min_i d_i\):

\[\deg u_i = (1,\, d_i - d_0), \qquad \deg x_j = (0,\, c_j), \qquad
\deg t = (1,\, -d_0),\]

so that \(u^a x^b\) has bidegree \((|a|,\, \deg(f^a x^b) - |a| d_0)\)
--- exactly what its image \(f^a x^b t^{|a|}\) records. The kernel is
bihomogeneous whatever the degrees of the \(f_i\), and when they agree
this \emph{is} the paper\textquotesingle s \(\deg u_i = (1,0)\).

\textbf{Normality reduces to \(S_2\).} A small projection satisfies
\(R_1\) automatically: a codimension-one point of \(Z\) cannot lie in
the exceptional locus, so \(\pi\) is a local isomorphism there and the
local ring is that of \(X\) at a codimension-one point, hence a DVR.
Testing smallness \emph{first} therefore reduces normality to \(S_2\).
This is not a small saving: on one example the \(R_1\) half of
\texttt{isNormal} took 119 s, passing through a Jacobian minor ideal
with 1 436 495 generators, against 0.02 s for \(S_2\).

\textbf{\(\omega_X\) embedded directly.} Steps 2 and 3 together amount
to embedding \(\omega_X\) into \(R\) as an ideal, so instead of choosing
\(s\) the implementation takes a homomorphism \(\omega_X \to R\) of
least degree. The \(s\) of Step 2 inherits whichever representative of
the canonical class the \texttt{Divisor} package returns; under the
grading \(\ell = (2,2,1)\) that gives \(s = y_3^{11}\) and an
\(I^{(1)}\) in degree 30, and the computation does not finish in
seventeen minutes. The direct embedding gives generators of degree 2 and
takes 0.05 s --- on the same threefold.

\textbf{Every divisor of \(\texttt{MaxSteps}!\).} One \(m\) that works
is as good as another, so what matters is finding a \emph{small} one,
and the cost per iteration grows steeply: on the threefold below, 0.02 s
at \(m=1\) against more than 600 s at \(m=12\). Trying every divisor of
\(\texttt{MaxSteps}!\) in increasing order still contains
\(1!, 2!, \dots\), so termination is unchanged, while filling the gaps
is nearly free.

\section{Environment and
reproducibility}\label{environment-and-reproducibility}

Computed with Macaulay2 1.24.11 on macOS (Apple Silicon), using the
bundled packages \texttt{Divisor}, \texttt{SymbolicPowers},
\texttt{Polyhedra}, \texttt{IntegralClosure}, \texttt{MinimalPrimes} and
\texttt{Elimination}. From the repository root:

\begin{verbatim}
M2 --script tests/run-tests.m2          # 13 regression tests, about 5 s
M2 --script examples/toric-flip.m2      # about 1 s
\end{verbatim}

Measured wall-clock times: the test suite 5 s, \texttt{toric-flip.m2}
1.0 s, \texttt{toric-flip-projective.m2} 1.3 s,
\texttt{toric-flip-index-two.m2} 2.7 s,
\texttt{ordinary-double-point.m2} 0.6 s. The checks quoted below ---
fans, signs, dimensions, smoothness, the agreement between code paths
--- are \texttt{assert}ed by the regression suite rather than merely
printed, so a regression shows up as a test failure. A few remarks are
derivations made here for the reader rather than assertions in the
suite, and are marked as such. The excerpts are trimmed; the complete
input and output are in the repository.

\section{Representative Macaulay2
calculations}\label{representative-macaulay2-calculations}

\subsection{A flip of a non-\QQ{}-Gorenstein toric
threefold}\label{a-flip-of-a-non-ux211a-gorenstein-toric-threefold}

Let \(\sigma \subset N_\RR{} = \RR^3\) be the cone on

\[v_1 = (1,0,0), \quad v_2 = (0,1,0), \quad v_3 = (0,0,1), \quad v_4 = (1,1,-2),\]

with circuit relation \(v_1 + v_2 = 2v_3 + v_4\), and let
\(X = \Spec\, k[\sigma^\vee \cap M]\). No \(m \in M_\QQ{}\) has
\(\langle m, v_i\rangle = 1\) for all four rays --- the first three
force \(m = (1,1,1)\) and then \(\langle m, v_4\rangle = 0\) --- so
\textbf{\(X\) is not \QQ{}-Gorenstein}, and this is the case the paper is
aimed at.

A circuit has exactly two triangulations, and toric intersection theory
gives the sign of \(K \cdot C\) on the wall curve from the wall relation
alone:

\begin{longtable}[]{@{}lll@{}}
\toprule\noalign{}
triangulation & wall & \(K \cdot C\) \\
\midrule\noalign{}
\endhead
\bottomrule\noalign{}
\endlastfoot
\(\langle v_1,v_2,v_3\rangle, \langle v_1,v_2,v_4\rangle\) &
\(\langle v_1,v_2\rangle\) & \(-(2+1-1-1) = -1\) \\
\(\langle v_1,v_3,v_4\rangle, \langle v_2,v_3,v_4\rangle\) &
\(\langle v_3,v_4\rangle\) & \(-(1+1-2-1) = +1\) \\
\end{longtable}

So the first is the flipping contraction \(Y \to X\) and \textbf{the
second is the flip \(Z \to X\)}. This is known before the algorithm is
run, and it is what the output is compared against. The algorithm
succeeds at \(m = 1\):

\begin{verbatim}
-- I = omega_X, embedded in R by a map of least degree
-- E has 2 component(s) with multiplicities {1, 1}
-- m = 1:
--   Z^m sits in P^1 x X, dim Z^m = 3
  component of dim 2 (not contracted)
  component of dim 2 (not contracted)
-- the flip is obtained for m = 1
\end{verbatim}

and the output is then dissected:

\begin{verbatim}
blown up ideal I = O_X(-E): ideal (y_5, y_4)
fibre over the fixed point has dimension 1
chart u_1 = 1:  dim 3,  smooth: true
chart u_2 = 1:  dim 3,  smooth: true

the fan of Z has maximal cones
  {{0, 1, 0}, {1, 1, -2}, {0, 0, 1}}
  {{1, 0, 0}, {1, 1, -2}, {0, 0, 1}}
equals the expected triangulation {v1,v3,v4}, {v2,v3,v4}:  true

Z: wall relation {1, 1, -1, -2},  K.C ~ 1
Y: wall relation {1, 2, -1, -1},  K.C ~ -1

K_Z is pi-ample and K_Y is pi-anti-ample, so Z is the flip of Y.
\end{verbatim}

Three things are confirmed at once. \(I\) has exactly two generators, so
\(Z\) lands in \(\PP^1 \times X\) --- one homogeneous coordinate per
maximal cone of \(Z\). The fan recovered from those generators is the
triangulation
\(\langle v_1,v_3,v_4\rangle, \langle v_2,v_3,v_4\rangle\). And the sign
of \(K \cdot C\), computed from the wall relation, is \(+1\): the
algorithm returned the flip and not the contraction it came from.

Also checked here: \texttt{canonicalDivisorData} returns \(-D_1 - D_3\)
rather than the torus-invariant \(-\sum_i D_i\), and the two differ by
the principal divisor \(\divisor(y_1) = D_2 + D_4\), so they are
linearly equivalent as they must be.

\subsection{The same flip over a projective
base}\label{the-same-flip-over-a-projective-base}

Every element of the Hilbert basis pairs to 1 with \(\ell = (1,1,-1)\),
so the standard grading is induced by a torus character and

\[X_{\mathrm{proj}} = \Proj\, k[y_1,\dots,y_5,w]/I \subset \PP^5,
\qquad \deg w = 1\]

is a projective threefold with \(D_+(w) = \Spec\, R\).
Projecting the rays of \(\tau = \sigma \times \RR_{\geq 0}\) shows its
only singularity is the vertex --- the point to be flipped --- so this
is the paper\textquotesingle s setting with nothing extra added.

\begin{verbatim}
degrees for l = (1,1,-1): {1, 1, 1, 1, 1}  (so the standard grading works)
singular locus is the vertex [0:...:0:1]: true
-- the flip is obtained for m = 1
I = ideal (y_5, y_4),  uniform degree used = null
affine I was ideal (y_5, y_4), the same ideal
Z_proj: 12 nonempty charts, dimensions {3}, all smooth: true
Z_proj restricted to D_+(w) equals the affine flip: true
\end{verbatim}

The last line is the check that matters: the graded and ungraded code
paths, which share very little, agree \emph{exactly} on the overlap.

\subsection{\texorpdfstring{A flip that needs
\(m = 2\)}{A flip that needs m = 2}}\label{a-flip-that-needs-m--2}

Replacing \(v_4\) by \((1,3,-2)\), with circuit relation
\(v_1 + 3v_2 = 2v_3 + v_4\), reverses which triangulation carries the
\(\pi\)-ample canonical class, and the flip becomes singular:
\(\langle v_1,v_2,v_4\rangle\) has determinant \(-2\), and
\(\langle m, v_1\rangle = \langle m, v_2\rangle = \langle m, v_4\rangle = 1\)
needs \(m = (1,1,3/2)\), so \(K_Z\) is 2-Cartier but not Cartier. This
is the first example in which the loop over \(m\) actually iterates:

\begin{verbatim}
-- m = 1:
--   Z^m sits in P^2 x X, dim Z^m = 3
  component of dim 2 (contracted)
  component of dim 2 (not contracted)
--   exceptional locus contains a divisor
-- m = 2:
--   Z^m sits in P^1 x X, dim Z^m = 3
  component of dim 2 (not contracted)
-- the flip is obtained for m = 2

I^(2) = ideal (y_6, y_4^2)
the fan of Z ... equals the expected triangulation {v1,v2,v3}, {v1,v2,v4}: true
chart u_1 = 1:  dim 3,  smooth: false
chart u_2 = 1:  dim 3,  smooth: true
\end{verbatim}

\(m = 1\) is rejected for the right reason --- a \emph{contracted}
component of dimension two, that is a divisor in the exceptional locus
--- and not by the normality test, which \(Z^1\) passes. The index of
\(K_Z\) being 2 and the required \(m\) being 2 agree, and one chart is
singular exactly as the determinant predicts.

\subsection{The ordinary double point, and the graph
morphism}\label{the-ordinary-double-point-and-the-graph-morphism}

The small resolution of the threefold ODP \(x_0x_1 = x_2x_3\) is the
standard sanity check, and it also exercises the Section 2 machinery:
the Segre construction of Lemma 2.3 turns the B2M projection into a
genuine morphism of monograded varieties (Lemma 2.6).

\begin{verbatim}
K_X = {{ideal(x2,x0), -3}, {ideal(x3,x0), -3}}
Z is normal:      true
pi is small:      true
-- Segre product has 10 generators
graph morphism with graph of dim 3
\end{verbatim}

The canonical divisor is an independent check, though this last reading
is a derivation rather than an assertion in the suite:
\(D_1 = V(x_0,x_2)\) and \(D_2 = V(x_0,x_3)\) are the two planes cut by
\(x_0 = 0\), so \(D_1 + D_2 = \divisor(x_0) \sim H\) and the
computed \(-3D_1 - 3D_2\) is \(-3H\) --- the right canonical class for a
quadric threefold.

\section{Why the answers are believed to be
correct}\label{why-the-answers-are-believed-to-be-correct}

No claim of formal verification is made. What the repository does
provide is a set of checks that are \textbf{independent of the algorithm
being tested}, in the sense that they are computed from the toric data
or from classical geometry rather than from the object the algorithm
returns.

\textbf{1. The answer is compared with an independently known answer.}
For a circuit there are exactly two triangulations, and which one is the
flip is determined by the sign of \(K \cdot C\) read off from the wall
relation. The tests recompute that sign from the fan and compare. The
algorithm is never asked to certify itself.

\textbf{2. Flip versus antiflip is checked, not assumed.} This is the
sharpest test, because both triangulations give a small birational
morphism to \(X\) and only the sign of \(K \cdot C\) separates them. A
sign error in Step 3 would return the contraction instead of its flip
and would otherwise look entirely healthy. This check is what gives
confidence in the sign the implementation uses: since
\(E = \divisor(s) - K_X\) is effective, the ideal \emph{of} \(\sO_X\)
corresponding to \(\sO_X(K_X - \divisor(s)) = \sO_X(-E)\) carries the
exponent \(m_i - n_i\), and the computation confirms that this choice
lands on the \(\pi\)-ample side.

\textbf{3. Independent code paths are made to agree.} The affine route
uses no grading at all; the projective route grades everything. They
share the divisor and Rees code but nothing of the degree bookkeeping.
On \(D_+(w)\) they produce literally the same ideal. Similarly, the two
routes to \(I\) --- the paper\textquotesingle s \(s\) and the direct
embedding of \(\omega_X\) --- produce the same divisor \(E\), and
\texttt{isS2Source} agrees with the full \texttt{isNormalSource} on
every example where the latter is affordable.

\textbf{4. Negative controls.} Tests that must \emph{fail} do fail:
blowing up a point of \(\PP^3\), or the origin of \(\bbA^3\), is divisorial
and is rejected by the smallness test; \(m = 1\) is rejected in the
index-two example. A test suite of positives only would not detect a
smallness test that always returns true.

\textbf{5. Classical invariants.} \(K_{\PP^3} = -4H\); the quadric
threefold has \(K = -3H\); the Segre product of two copies of \(\PP^1\)
has four generators and \(\PP^1 \times \PP{}(1,2)\) has five. These pin down
the supporting routines independently of any flip.

\textbf{6. The checks have demonstrably caught a real error.} This is
the strongest evidence that they are not vacuous. Levelling the
generators of \(I^{(m)}\) to a single degree --- the natural way to make
the presentation in \(X \times \PP^{r-1}\) work, and what this package did
originally --- reports the toric flip above as \textbf{not small} when
the base is graded by \(\ell = (3,2,1)\), while the affine route reports
it as small. Seven of eight mixed-degree cases tried had the wrong
saturation. The discrepancy between two routes that must agree is what
exposed it; it is now a regression test, and the corrected grading is
described above.

\subsection{What this does not
establish}\label{what-this-does-not-establish}

\begin{itemize}
\tightlist
\item
  It is not a proof of the algorithm, and not a proof that the
  implementation matches the paper. It is evidence that on these inputs
  the two agree.
\item
  All flips computed here are \textbf{toric}, and all are settled by
  \(m = 1\) or \(m = 2\). Nothing is known about the behaviour at large
  \(m\) beyond the timings.
\item
  Over a weighted grading the \(S_2\) test is \emph{sufficient} but not
  proved necessary, so the loop may in principle try one more \(m\) than
  it has to. This affects optimality, not correctness.
\item
  No check is made that the input really is the target of a flipping
  contraction. The algorithm assumes it, and will return something for
  input that is not.
\item
  The base field is \QQ{} in every example; the paper works over
  \(\overline{\QQ{}}\). Nothing in the code assumes \QQ{}, but number fields are
  untested.
\end{itemize}

\section{Summary of results}\label{summary-of-results}

\begin{longtable}[]{@{}llccp{0.44\linewidth}@{}}
\toprule\noalign{}
Example & Base & \(m\) & \(\dim Z\) & Independent check \\
\midrule\noalign{}
\endhead
\bottomrule\noalign{}
\endlastfoot
Toric \((1,1,-2)\) & \(\Spec\) & 1 & 3 & fan \(=\) expected
triangulation; \(K \cdot C = +1\) \\
Toric \((1,1,-2)\) & \(\Proj\) & 1 & 3 & restricts to the affine
answer on \(D_+(w)\); 12 smooth charts \\
Toric \((1,3,-2)\) & \(\Spec\) & 2 & 3 & \(m=1\) rejected;
\(K_Z\) of index 2; one chart singular \\
Toric, \(\ell=(3,2,1)\) & \(\Proj\) & 1 & 3 & agrees with
affine; levelling gives the wrong answer \\
ODP, small resolution & both & --- & 3 & normal and small; \(K = -3H\);
Segre model of dimension 3 \\
\(\Bl_{pt} \PP^3\), \(\Bl_0 \bbA^3\) & both & --- & 3 &
correctly rejected as not small \\
\end{longtable}

\textbf{Interpretation.} On these inputs the implementation reproduces
flips that are known independently, distinguishes them from their
antiflips, rejects divisorial contractions, and agrees with itself
across two code paths and two constructions of \(I\). That is evidence
that it behaves as intended here. It is neither a proof of the algorithm
nor a statement about inputs of a different shape.

\section{Limitations}\label{limitations}

\begin{itemize}
\tightlist
\item
  Only toric examples have been run end to end, and only in dimension
  three.
\item
  Only \(m \leq 2\) has been needed; larger \(m\) is exercised for cost,
  not for correctness.
\item
  The Segre construction of Lemma 2.3 requires the fibre variables to
  have bidegree \((d_j, 0)\), so \texttt{b2mToGraphMorphism} does not
  apply once the fibre weights are nontrivial, and raises an error.
  Generalising it is the main open problem left by the correction
  described above.
\item
  There is no affine analogue of the graph morphism at all.
\item
  \texttt{isS2Source} dominates the cost at large \(m\) (5.9 s of 7.0 s
  at \(m = 8\)).
\item
  Number fields, and bases that are not \QQ{}, are untested.
\end{itemize}

\section{Use of AI in the
implementation}\label{use-of-ai-in-the-implementation}

The Macaulay2 code, the tests, and this note were written essentially by
an AI system (Claude), with edits by the author. The author has read
them and believes them correct, but has not checked every line or every
intermediate computation.

Two things nevertheless make large errors unlikely. First, the algorithm
is the one in the paper below; little is attempted here that is
conceptually new, so there is limited room for the code to be quietly
doing something else. Second, and more to the point, the tests do not
merely record what the code computes: each example is checked against
geometry established independently --- the triangulations of a circuit,
the sign of \(K \cdot C\) from a wall relation, determinants of cones,
canonical classes of $\PP^3$ and of a quadric threefold --- and those checks
run automatically. A serious error would have to reproduce all of them
by accident.

That this is not an idle argument is shown by the levelling bug in item
6 above: an error that had passed every existing test was found
precisely because two routes that must agree did not.

None of this replaces a proof.

\section{Conclusion}\label{conclusion}

Algorithm 3 runs, and it produces flips of non-\QQ{}-Gorenstein toric
threefolds that can be checked against toric geometry --- including the
case where the flip is singular and the multiplier \(m = 1\) genuinely
fails. The most useful outcome is not the number of examples but the
chain connecting the construction in the paper, the source code, the
actual input, and output that is verified against independently known
answers. Along the way the implementation turned up one point --- the
presentation of \(Z^m\) in \(X \times \PP^{r-1}\), and the levelling of
generators it appears to require --- that needs a hypothesis or a
reformulation once the base carries a weighted grading; it is recorded
above and in the implementation notes.

\section{References}\label{references}

\begin{itemize}
\item
  \textbf{Yasuda, T.} (2026). An algorithm for the minimal model program
  in dimension three. arXiv:2603.13703.
  \url{https://arxiv.org/abs/2603.13703} (Section 2 and Algorithm 3 of
  Section 7; this repository implements Algorithm 3.)
\item
  \textbf{Cox, D. A., Little, J. B. \& Schenck, H. K.} (2011).
  \emph{Toric Varieties.} Graduate Studies in Mathematics 124, American
  Mathematical Society. (Circuits, wall relations, and the sign of
  \(K \cdot C\) used in the checks.)
\item
  \textbf{Kollár, J. \& Mori, S.} (1998). \emph{Birational Geometry of
  Algebraic Varieties.} Cambridge Tracts in Mathematics 134. (Flips and
  flipping contractions.)
\item
  \textbf{Grayson, D. R. \& Stillman, M. E.} Macaulay2, a software
  system for research in algebraic geometry.
  \url{https://macaulay2.com/}
\item
  \textbf{Schwede, K., Smith, Z. et al.} The \texttt{Divisor} package
  for Macaulay2. (Canonical divisors and divisorial ideals; used in
  place of \texttt{SymbolicPowers}.)
\end{itemize}

\end{document}
